\documentclass[12pt]{article}
\usepackage[utf8]{inputenc}
\usepackage[margin=1in]{geometry}
\usepackage{times}

\begin{document}

\begin{flushleft}
	\textbf{Name:} Harshit Prashant Dhanwalkar \\
	\textbf{Branch:} Earth System Sciences (ESS)\\
	\textbf{SCcode:} SC21B164\\
\end{flushleft}

\vspace{1cm}

\begin{center}
	\Large \textbf{Abstract} \\
	\vspace{0.5cm}
	\large \textbf{Impact of Land-Surface Feedback on the Post-Landfall Sustainability of Monsoon Depressions}
\end{center}

\vspace{0.5cm}

\noindent Investigation on the mechanisms governing the post-landfall sustainability of Monsoon Depressions (MDs) over the Indian subcontinent during the 2024 monsoon season, with a specific focus on the role of soil moisture (Soil water Volumatric Level-1, \texttt{SWVL1}) and total precipitation (\texttt{TP}). Utilizing ERA5 reanalysis and high-resolution tracking algorithms, analysing the thermodynamic transitions occurring at the land-sea interface. The Brown Ocean Effect, antecedent soil moisture and ongoing precipitation serve as a latent heat reservoir, potentially offsetting the typical decay associated with inland movement.

\vspace{0.3cm}

\noindent Azimuthal Hovm\"{o}ller analysis is employed to visualize the radial evolution of key variables relative to the MD center over time. Preliminary findings indicate a strong positive correlation between land-phase \texttt{TP} and \texttt{SWVL1} (layer-1, 0-7cm), suggesting a feedback loop. Integrating correlation significance with soil water capacity to distinguish between systems sustained by local surface forcing versus those driven by large-scale dynamics. The results provide insights into why certain 2024 MDs exhibited extended inland longevity, highlighting the critical role of saturated soil conditions in maintaining system intensity after landfall.

\end{document}















